
Gradient attribution methods fail to detect unknown spurious correlations --- practitioners cannot identify which features drive model shortcuts using saliency maps alone. We extend GAIA (Gradient Abnormality for Out-of-Distribution Detection) from distribution shift to subpopulation shift, proposing a gradient-based framework for detecting and mitigating spurious correlations via gradient process abnormality rather than attribution results. The methodology comprises three stages: (1) GradCAM extraction, (2) GAIA-Z zero-deflation computation, and (3) statistical testing. We validate the approach through synthetic experiments and proof-of-concept tests. Synthetic validation (5794 samples) demonstrates the pipeline differentiates controlled gradient patterns (GAIA-Z divergence 0.30, p<0.0001, Cohen's d=198.75). Causality tests show synthetic background swap reduces abnormality by 39.4\% (p<0.001, d=2.96). Spatial gradient regularization methodology improves worst-group accuracy by 23 percentage points on MNIST+Color (78\% vs 55\%, 1 seed, 2 epochs). All validation results are synthetic or minimal proof-of-concept. Real Waterbirds experiments require GPU compatibility resolution (PyTorch 2.1+ for H100 sm\_90 support) or CPU training (estimated 70-90 hours). We position this as a Tier 3 methodology contribution with empirical validation pending.
